\section{Capítulo~\ref{sec:dendrites}. Número de dendritos}

A estrutura das classes e o número de entradas foram medidos por anatomia e por registros elétricos; a estimativa~\eqref{eq:dendriteload} é física simples do capacitor, e a explicação computacional do número de entradas se apoia em teoria e modelos.

{\footnotesize
\setlength{\tabcolsep}{3pt}\renewcommand{\arraystretch}{1.15}
\begin{longtable}{>{\raggedright}p{0.5cm}>{\raggedright}p{4.0cm}>{\raggedright}p{5.6cm}>{\raggedright}p{2.3cm}>{\raggedright\arraybackslash}p{1.7cm}}
\toprule
N.º & Proposição & Experimento e o que foi medido & Fonte & Status \\
\midrule\endfirsthead
\toprule
N.º & Proposição & Experimento e o que foi medido & Fonte & Status \\
\midrule\endhead
1 & Capacitância específica da membrana $c_m\approx1$~\textmu F/cm$^2$ & Retirava-se membrana do corpo do neurônio com uma pipeta, aplicava-se um degrau de tensão e, pelo decaimento da corrente capacitiva, achava-se a capacitância total, depois dividida pela área: $0.9$~\textmu F/cm$^2$ em três classes de neurônios & \cite{Gentet2000} & medido \\
2 & Tempo de carga $t\approx c_mA\Delta V/I$~\eqref{eq:dendriteload}: o neurônio grande precisa de dezenas de entradas simultâneas, ao pequeno basta uma sinapse grande & Estimativa pela lei de carga do capacitor; os números ($20$ e $300$~pF, $0.1$ e alguns nA) são ordens de grandeza & dedução no capítulo & teoria \\
3 & Uma sinapse enorme dá corrente de alguns nanoampères e leva o neurônio pequeno ao limiar de imediato & Registro simultâneo do terminal e do destinatário numa fatia: corrente de uma sinapse de $2$--$13$~nA, cada disparo de entrada provoca resposta supralimiar, atraso de cerca de $0.4$~ms a $36^\circ$C; o destinatário emite \nt{GLYC} & \cite{Borst1995,BorstSoria2012} & medido \\
4 & Zero dendritos: nos neurônios sensoriais do tato e da dor, o disparo vai do sensor à medula espinhal sem passar pelo corpo celular & A estrutura (axônio que se divide em dois junto ao corpo) foi estabelecida anatomicamente. Que a condução não exige excitabilidade do corpo foi mostrado num modelo validado desse neurônio: sem corpo excitável, o disparo atravessa a bifurcação com a mesma confiabilidade & \cite{AmirDevor2003} & teoria \\
5 & Um dendrito: o neurônio olfatório carrega um tipo de receptor e transmite uma linha de entrada & Revisão de experimentos com genes de receptores: cada neurônio olfatório expressa um gene de receptor de uma grande família & \cite{Mombaerts2004} & medido \\
6 & Dois dendritos: nos detectores de direção do som, as entradas do ouvido esquerdo e do direito chegam a dendritos diferentes & Anatomia e registros em mamíferos, reunidos numa revisão: as entradas dos dois ouvidos são separadas em dois dendritos opostos & \cite{Grothe2010} & medido \\
7 & Separar as entradas em dois dendritos torna o ``E'' mais estrito & Mostrado em modelo: a saturação~\eqref{eq:shunt} num só dendrito impede que as entradas de um lado cheguem ao limiar, e a detecção de coincidências melhora. Não há experimento direto em que se tenha mudado a distribuição das entradas & \cite{AgmonSnir1998} & teoria \\
8 & Cerca de $7\cdot10^{10}$ pequenos neurônios \nt{GLUT} do circuito de aprendizado motor, com $\sim4$ entradas cada & Contagem de células por fracionamento isotrópico: nessa parte do cérebro humano há cerca de $69$ bilhões de neurônios. Registro de uma dessas células em rato vivo: ao toque chegam rajadas de correntes discretas de poucas entradas, e a célula dispara quando elas se somam & \cite{Azevedo2009,Chadderton2004} & medido \\
9 & Um elemento de limiar com $n$ entradas separa no máximo $P_{\max}\approx2n$ padrões aleatórios~\eqref{eq:cover}; para o neurônio de saída do circuito de aprendizado motor, o limite é $\sim2\cdot10^5$, e a estimativa realista, $\sim4\cdot10^4$ associações & O limite é um resultado clássico da geometria combinatória. A estimativa realista vem da teoria de capacidade com pesos só positivos e margem para o ruído, aplicada ao número medido de entradas desse neurônio & \cite{Cover1965,Brunel2004} & teoria \\
10 & Os misturadores de poucas entradas servem para expandir a entrada; ``quatro'' está perto do ótimo & Teoria: a dimensão da representação dada pela camada de misturadores é máxima aproximadamente com o número de entradas observado na anatomia; a hipótese original da expansão está nos trabalhos clássicos & \cite{Marr1969,Albus1971,LitwinKumar2017} & hipótese \\
11 & Árvores grandes: $10^4$--$10^5$ entradas & Contagem de sinapses em imagens de microscopia eletrônica: $\approx1.7\cdot10^5$ sinapses num neurônio \nt{GABA} de saída do circuito de aprendizado motor em rato; cerca de $30\,000$ entradas excitatórias e $1\,700$ inibitórias num neurônio \nt{GLUT} do hipocampo & \cite{NapperHarvey1988,Megias2001} & medido \\
12 & Os ramos de uma árvore grande são subcircuitos separados com limiares próprios; o neurônio é uma pequena rede de várias camadas & Estimulação pontual de dois locais em dendritos finos: entradas no mesmo ramo se somam de forma sigmoide, em ramos diferentes, linearmente. Nos dendritos de neurônios do córtex humano foram encontrados disparos graduados de cálcio que permitem a uma só célula resolver uma tarefa linearmente não separável. Modelo: para reproduzir um neurônio é preciso uma rede profunda & \cite{Polsky2004,Gidon2020,Beniaguev2021} & medido \\
13 & Dendritos-saída: cada ramo do neurônio \nt{GABA} da retina é um detector de direção separado & Registro de cálcio por dois fótons em ramos isolados: o sinal no ramo é seletivo ao movimento do corpo celular para a ponta, e a tensão do corpo não & \cite{Euler2002} & medido \\
14 & O número de entradas segue a tarefa (proposição~\ref{prop:dendrites}) & A estrutura das classes foi medida (linhas 3--13); que o número de entradas foi escolhido pela velocidade ou pela classificação só foi mostrado por teoria e modelos, para algumas classes & \cite{LitwinKumar2017,AgmonSnir1998} & hipótese \\
\bottomrule
\end{longtable}}
